\section{2016级工科数学分析下 A卷}

\subsection{资料来源}
\begin{itemize}
  \item 试题：2016 级工科数学分析下 A 卷，已导出页面图校正公式。
  \item 未发现配套 A 卷答案或解析文本；下列答案与解析为据题面补写。
\end{itemize}

\subsection{试题}

\paragraph*{试卷信息}
诚信应考，考试作弊将带来严重后果！

华南理工大学本科生期末考试《工科数学分析（二）》2016--2017学年第二学期期末考试试卷 A 卷。

注意事项：开考前请将密封线内各项信息填写清楚；所有答案请直接答在试卷上；考试形式：闭卷；本试卷共5大题，满分100分，考试时间120分钟。

\paragraph*{一、填空题（共5题，每题2分，共10分）}
\begin{enumerate}
  \item 函数
  \[
    f(x,y)=2x^2+y^2
  \]
  在点 \((1,1)\) 处沿该点的梯度方向的方向导数为 \underline{\hspace{5cm}}。
  \item 向量场 \((2x-3y)\mathbf i+(3x-z)\mathbf j+(y-2x)\mathbf k\) 的旋度向量为 \underline{\hspace{5cm}}。
  \item 设 \(\Sigma\) 表示球面 \(x^2+y^2+z^2=R^2\) 的外侧，则第二类曲面积分
  \[
    \iint_\Sigma x\,\mathrm{d}y\mathrm{d}z+y\,\mathrm{d}z\mathrm{d}x+z\,\mathrm{d}x\mathrm{d}y=
  \]
  \underline{\hspace{5cm}}。
  \item 初值问题
  \[
    \begin{cases}
      y'+y\cos x=\mathrm{e}^{-\sin x},\\
      y|_{x=0}=1
    \end{cases}
  \]
  的解为 \underline{\hspace{5cm}}。
  \item 设 \(f(x)\) 是周期为 \(2\) 的周期函数，它在 \((-1,1]\) 上的表达式为
  \[
    f(x)=
    \begin{cases}
      2, & -1<x\le 0,\\
      x^3+1, & 0<x\le 1,
    \end{cases}
  \]
  则 \(f(x)\) 的傅里叶级数在 \(x=0\) 处收敛于 \underline{\hspace{5cm}}。
\end{enumerate}

\paragraph*{二、单选题（每题只有一个正确选项，共5题，每题2分，共10分）}
\begin{enumerate}
  \item 二元函数
  \[
    f(x,y)=
    \begin{cases}
      \dfrac{xy}{x^2+y^2}, & (x,y)\ne(0,0),\\
      0, & (x,y)=(0,0)
    \end{cases}
  \]
  在点 \((0,0)\) 处（\quad）。
  \begin{enumerate}
    \item[A.] 连续，偏导数存在。
    \item[B.] 不连续，偏导数存在。
    \item[C.] 连续，偏导数不存在。
    \item[D.] 不连续，偏导数不存在。
  \end{enumerate}

  \item 曲线
  \[
    \begin{cases}
      x^2+y^2+z^2=9,\\
      x-y+z=1
    \end{cases}
  \]
  在点 \(M(1,2,2)\) 处的切线一定平行于（\quad）。
  \begin{enumerate}
    \item[A.] \(Oxy\) 平面。
    \item[B.] \(Oyz\) 平面。
    \item[C.] \(Ozx\) 平面。
    \item[D.] 平面 \(x-y+z=1\)。
  \end{enumerate}

  \item 设 \(D\) 是一个有界的平面闭区域，其边界曲线 \(\Gamma\) 分段光滑，则下列积分值不等于区域 \(D\) 的面积的是（\quad）。
  \begin{enumerate}
    \item[A.] \(\displaystyle \oint_\Gamma y\,\mathrm{d}x\)
    \item[B.] \(\displaystyle \frac12\oint_\Gamma x\,\mathrm{d}y-y\,\mathrm{d}x\)
    \item[C.] \(\displaystyle \oint_\Gamma x\,\mathrm{d}y\)
    \item[D.] \(\displaystyle \iint_D 1\,\mathrm{d}x\mathrm{d}y\)
  \end{enumerate}

  \item 关于未知函数 \(y\) 的微分方程 \((y-\sin x)\,\mathrm{d}x+x\,\mathrm{d}y=0\) 是（\quad）。
  \begin{enumerate}
    \item[A.] 可分离变量方程。
    \item[B.] 一阶非齐次线性方程。
    \item[C.] 一阶齐次线性方程。
    \item[D.] 非线性方程。
  \end{enumerate}

  \item 使得级数 \(\displaystyle \sum_{n=1}^{\infty}\frac{(-1)^n}{n^p}\) 条件收敛的常数 \(p\) 的取值范围是（\quad）。
  \begin{enumerate}
    \item[A.] \(p\le 0\)
    \item[B.] \(0<p<1\)
    \item[C.] \(0<p\le 1\)
    \item[D.] \(p>1\)
  \end{enumerate}
\end{enumerate}

\paragraph*{三、计算题（共4题，每题7分，共28分）}
\begin{enumerate}
  \item 设函数 \(u=f(x,y,z)\) 有连续的偏导数，函数 \(y=y(x)\) 和 \(z=z(x)\) 分别由方程
  \(\mathrm{e}^{xy}=y\) 和 \(\mathrm{e}^z=xz\) 确定，计算 \(\dfrac{\mathrm{d}u}{\mathrm{d}x}\)。

  \item 计算累次积分
  \[
    \int_1^2\mathrm{d}x\int_{\sqrt{x}}^x\sin\frac{\pi x}{2y}\,\mathrm{d}y
    +\int_2^4\mathrm{d}x\int_{\sqrt{x}}^2\sin\frac{\pi x}{2y}\,\mathrm{d}y.
  \]

  \item 计算三重积分 \(\displaystyle \iiint_\Omega z\,\mathrm{d}x\mathrm{d}y\mathrm{d}z\)，其中 \(\Omega\) 是由上半球面
  \(x^2+y^2+z^2=R^2\)，\(z\ge 0\)，\(R>0\) 和锥面 \(z=\sqrt{x^2+y^2}\) 所围成的闭区域。

  \item 计算第二类曲线积分
  \[
    \oint_\Gamma xy^2\,\mathrm{d}y-x^2y\,\mathrm{d}x,
  \]
  其中 \(\Gamma\) 为椭圆 \(\dfrac{x^2}{4}+y^2=1\)，取逆时针方向。
\end{enumerate}

\paragraph*{四、解答题（共4题，每题8分，共32分）}
\begin{enumerate}
  \item 设二阶常系数线性微分方程
  \[
    y''+\alpha y'+\beta y=-\mathrm{e}^x
  \]
  的一个特解为 \(y=(1+x)\mathrm{e}^x\)，试确定常数 \(\alpha,\beta\)，并求该方程的通解。

  \item 已知螺旋形弹簧一圈的方程为
  \[
    x=a\cos t,\qquad y=a\sin t,\qquad z=bt,\qquad 0\le t\le 2\pi,
  \]
  其中 \(a,b\) 为大于零的常数，且弹簧上各点处的线密度等于该点到 \(Oxy\) 平面的距离，求此弹簧的质心坐标。

  \item 求上半球面 \(x^2+y^2+z^2=a^2,\ z\ge 0\) 被柱面 \(x^2+y^2=ax\) 截下的部分的面积。

  \item 将函数 \(f(x)=\arctan x\) 展开为 \(x\) 的幂级数。
\end{enumerate}

\paragraph*{五、证明题（共2题，每题6分，共12分）}
\begin{enumerate}
  \item 设函数 \(f(\xi,\eta)\) 具有连续的二阶偏导数，且满足
  \[
    \frac{\partial^2f}{\partial \xi^2}+\frac{\partial^2f}{\partial \eta^2}=0.
  \]
  证明：函数 \(z=f(x^2-y^2,2xy)\) 满足
  \[
    \frac{\partial^2z}{\partial x^2}+\frac{\partial^2z}{\partial y^2}=0.
  \]

  \item 证明：函数项级数 \(\displaystyle \sum_{n=0}^{\infty}(1-x)x^n\) 在区间 \((0,1)\) 上点态收敛，但不一致收敛。
\end{enumerate}

\paragraph*{六、应用题（共1题，共8分）}
求曲线
\[
  \begin{cases}
    z=x^2+y^2,\\
    y=\dfrac1x
  \end{cases}
\]
上到 \(Oxy\) 平面距离最近的点。

\subsection{答案与解析}
\paragraph*{一、填空题}
\begin{enumerate}
  \item \(\nabla f(1,1)=(4,2)\)，沿梯度方向的方向导数为 \(\|\nabla f(1,1)\|=2\sqrt5\)。
  \item
  \[
    \nabla\times\mathbf F=
    \begin{vmatrix}
      \mathbf i&\mathbf j&\mathbf k\\
      \partial_x&\partial_y&\partial_z\\
      2x-3y&3x-z&y-2x
    \end{vmatrix}
    =2\mathbf i+2\mathbf j+6\mathbf k.
  \]
  \item 由 Gauss 公式，
  \[
    \iint_\Sigma x\,\mathrm{d}y\mathrm{d}z+y\,\mathrm{d}z\mathrm{d}x+z\,\mathrm{d}x\mathrm{d}y
    =\iiint_V3\,\mathrm{d}V=4\pi R^3.
  \]
  \item 积分因子为 \(\mathrm{e}^{\sin x}\)，故
  \[
    (y\mathrm{e}^{\sin x})'=1,\qquad y=(x+1)\mathrm{e}^{-\sin x}.
  \]
  \item 傅里叶级数在跳跃点收敛于左右极限平均值：
  \[
    \frac{f(0-0)+f(0+0)}2=\frac{2+1}{2}=\frac32.
  \]
\end{enumerate}

\paragraph*{二、单选题}
\begin{enumerate}
  \item 选 B。沿 \(y=x\) 趋于原点时 \(f(x,x)=1/2\)，故不连续；沿坐标轴计算偏导数均存在。
  \item 选 D。切向量垂直于两个曲面的法向量
  \[
    (2,4,4)\times(1,-1,1)=(8,2,-6),
  \]
  它与平面 \(x-y+z=1\) 的法向量 \((1,-1,1)\) 垂直，故切线平行于该平面。
  \item 选 A。正向边界下 \(\oint_\Gamma y\,\mathrm{d}x=-S_D\)，而其余三项均等于区域面积。
  \item 选 B。方程化为
  \[
    y'+\frac1x y=\frac{\sin x}{x},
  \]
  是一阶非齐次线性方程。
  \item 选 C。交错级数收敛要求 \(p>0\)，绝对收敛要求 \(p>1\)，故条件收敛为 \(0<p\le 1\)。
\end{enumerate}

\paragraph*{三、计算题}
\begin{enumerate}
  \item 由 \(\mathrm{e}^{xy}=y\) 得
  \[
    y'(1-xy)=y^2,\qquad y'=\frac{y^2}{1-xy}.
  \]
  由 \(\mathrm{e}^z=xz\) 得
  \[
    z'(\mathrm{e}^z-x)=z,\qquad z'=\frac{z}{x(z-1)}.
  \]
  因此
  \[
    \frac{\mathrm{d}u}{\mathrm{d}x}
    =f_x+f_y\frac{y^2}{1-xy}+f_z\frac{z}{x(z-1)}.
  \]

  \item 积分区域合并后为
  \[
    1\le y\le 2,\qquad y\le x\le y^2.
  \]
  所以
  \[
    I=\int_1^2\mathrm{d}y\int_y^{y^2}\sin\frac{\pi x}{2y}\,\mathrm{d}x
    =-\frac2\pi\int_1^2 y\cos\frac{\pi y}{2}\,\mathrm{d}y
    =\frac4{\pi^2}+\frac8{\pi^3}.
  \]

  \item 用球坐标，区域为
  \[
    0\le r\le R,\qquad 0\le\varphi\le\frac{\pi}{4},\qquad 0\le\theta\le2\pi.
  \]
  故
  \[
    \iiint_\Omega z\,\mathrm{d}V
    =\int_0^{2\pi}\!\!\int_0^{\pi/4}\!\!\int_0^R r\cos\varphi\cdot r^2\sin\varphi\,\mathrm{d}r\,\mathrm{d}\varphi\,\mathrm{d}\theta
    =\frac{\pi R^4}{8}.
  \]

  \item 令 \(P=-x^2y\)，\(Q=xy^2\)。由 Green 公式，
  \[
    \oint_\Gamma P\,\mathrm{d}x+Q\,\mathrm{d}y
    =\iint_D(x^2+y^2)\,\mathrm{d}A.
  \]
  对椭圆作变换 \(x=2r\cos\theta,\ y=r\sin\theta\)，得
  \[
    \iint_D(x^2+y^2)\,\mathrm{d}A
    =\int_0^{2\pi}\!\!\int_0^1(4r^2\cos^2\theta+r^2\sin^2\theta)\,2r\,\mathrm{d}r\,\mathrm{d}\theta
    =\frac{5\pi}{2}.
  \]
\end{enumerate}

\paragraph*{四、解答题}
\begin{enumerate}
  \item 将 \(y=(1+x)\mathrm{e}^x\) 代入方程，得
  \[
    \bigl[(1+\alpha+\beta)x+(3+2\alpha+\beta)\bigr]\mathrm{e}^x=-\mathrm{e}^x.
  \]
  所以
  \[
    1+\alpha+\beta=0,\qquad 3+2\alpha+\beta=-1,
  \]
  解得 \(\alpha=-3,\ \beta=2\)。方程为
  \[
    y''-3y'+2y=-\mathrm{e}^x,
  \]
  齐次解为 \(C_1\mathrm{e}^x+C_2\mathrm{e}^{2x}\)，故通解可写为
  \[
    y=C_1\mathrm{e}^x+C_2\mathrm{e}^{2x}+x\mathrm{e}^x.
  \]

  \item 有
  \[
    \mathrm{d}s=\sqrt{a^2+b^2}\,\mathrm{d}t,\qquad \rho=bt.
  \]
  质量
  \[
    m=\int_0^{2\pi}bt\sqrt{a^2+b^2}\,\mathrm{d}t
    =2\pi^2b\sqrt{a^2+b^2}.
  \]
  因而
  \[
    \bar x=0,\qquad
    \bar y=\frac{\int_0^{2\pi}a\sin t\cdot bt\sqrt{a^2+b^2}\,\mathrm{d}t}{m}
    =-\frac a\pi,
  \]
  \[
    \bar z=\frac{\int_0^{2\pi}bt\cdot bt\sqrt{a^2+b^2}\,\mathrm{d}t}{m}
    =\frac{4\pi b}{3}.
  \]
  所以质心为
  \[
    \left(0,-\frac a\pi,\frac{4\pi b}{3}\right).
  \]

  \item 投影区域为柱面内部
  \[
    D:\ 0\le r\le a\cos\theta,\qquad -\frac{\pi}{2}\le\theta\le\frac{\pi}{2}.
  \]
  上半球面 \(z=\sqrt{a^2-r^2}\) 上
  \[
    \mathrm{d}S=\frac{a}{\sqrt{a^2-r^2}}r\,\mathrm{d}r\,\mathrm{d}\theta.
  \]
  故面积
  \[
    S=\int_{-\pi/2}^{\pi/2}\int_0^{a\cos\theta}
    \frac{ar}{\sqrt{a^2-r^2}}\,\mathrm{d}r\,\mathrm{d}\theta
    =a^2(\pi-2).
  \]

  \item 由
  \[
    \frac1{1+x^2}=\sum_{n=0}^{\infty}(-1)^n x^{2n}\qquad (|x|<1)
  \]
  逐项积分得
  \[
    \arctan x=\sum_{n=0}^{\infty}\frac{(-1)^n x^{2n+1}}{2n+1},
    \qquad |x|\le 1.
  \]
\end{enumerate}

\paragraph*{五、证明题}
\begin{enumerate}
  \item 记 \(\xi=x^2-y^2,\ \eta=2xy\)。链式法则给出
  \[
    \frac{\partial^2z}{\partial x^2}+\frac{\partial^2z}{\partial y^2}
    =4(x^2+y^2)\left(\frac{\partial^2f}{\partial \xi^2}+\frac{\partial^2f}{\partial \eta^2}\right)=0.
  \]

  \item 对固定 \(x\in(0,1)\)，
  \[
    \sum_{n=0}^{\infty}(1-x)x^n=1,
  \]
  因而点态收敛。第 \(N\) 项部分和为
  \[
    S_N(x)=1-x^{N+1},
  \]
  余项为 \(x^{N+1}\)。由于
  \[
    \sup_{0<x<1}|1-S_N(x)|=\sup_{0<x<1}x^{N+1}=1,
  \]
  余项不一致趋于 \(0\)，故不一致收敛。
\end{enumerate}

\paragraph*{六、应用题}
曲线上点到 \(Oxy\) 平面的距离为 \(|z|\)。由题设
\[
  z=x^2+\frac1{x^2}\qquad (x\ne0),
\]
故只需求 \(z\) 的最小值。令
\[
  z'=2x-\frac2{x^3}=0,
\]
得 \(x^4=1\)，即 \(x=\pm1\)。此时 \(y=1/x\)，\(z=2\)。所以最近点为
\[
  (1,1,2),\qquad (-1,-1,2).
\]
